% SPDX-License-Identifier: GPL-3.0-or-later
% Copyright (C) 2026 Alessandro Greco
\chapter{Публикация onion-сервиса}\label{cap:21-onion-pratica}
\aggiornatoal{23 сентября 2026 г.}{tor 0.4.9.13; документация Onionbalance}

\begin{obiettivi}
К концу этой главы вы сможете:
\begin{itemize}[nosep]
  \item настраивать onion-сервис и защищать его ключи;
  \item избегать ошибок конфигурации, раскрывающих сервер;
  \item предлагать посетителям обычного веб-сайта его версию \texttt{.onion} через Onion-Location;
  \item ограничивать доступ авторизованными клиентами с помощью ограниченного обнаружения (restricted discovery, Лабораторная работа~13);
  \item оценивать адреса с узнаваемым префиксом (vanity), одношаговые сервисы и балансировку нагрузки через Onionbalance.
\end{itemize}
\end{obiettivi}

В \cref{cap:14-onion-services} мы разобрали, как работают onion-сервисы. Эта глава — практическое продолжение: публикация сервиса, предотвращение простых ошибок, способных нарушить его анонимность, и выбор инструментов для типичных эксплуатационных задач.

% ---------------------------------------------------------------------------
\section{Первый onion-сервис}\label{sec:onionp-primo}

Достаточно двух строк в \file{torrc}:
\begin{lstlisting}[language=torrc]
HiddenServiceDir /var/lib/tor/my_service/
HiddenServicePort 80 127.0.0.1:8080
\end{lstlisting}
Путь \opt{HiddenServiceDir} должен быть доступен процессу tor для чтения и записи; в лаборатории это путь \emph{внутри контейнера Chutney}. В первом примере создаётся \emph{общедоступный} onion-сервис, которому файл \file{alice.auth} не нужен. Авторизация клиентов и файлы \file{authorized\_clients/*.auth} рассматриваются далее в этой главе.
Первая задаёт каталог, в котором tor хранит ключи сервиса. Вторая сопоставляет \emph{виртуальный} порт 80 его адреса \texttt{.onion} с локальным сервером на порту 8080. После перезапуска tor создаёт в этом каталоге следующие файлы:
\begin{description}
  \item[\file{hostname}] Адрес \texttt{.onion} (\cref{sec:os-indirizzo}).
  \item[\file{hs\_ed25519\_public\_key}, \file{hs\_ed25519\_secret\_key}] Пара идентификационных ключей сервиса. \textbf{Тот, кто владеет секретным ключом, контролирует адрес}: он может публиковать дескрипторы и выдавать себя за сервис. Защищайте этот ключ так же, как закрытый ключ сертификата, и храните безопасную резервную копию. Если ключ потерян, адрес утрачен навсегда.
\end{description}
Каталог должен быть доступен только пользователю, от имени которого работает tor; в противном случае tor откажется его использовать. Один процесс tor может размещать несколько сервисов, каждый в своём каталоге. В лаборатории мы добавили второй сервис к \texttt{test009h} ещё двумя строками.

\begin{figure}[htbp]
\centering
% SPDX-License-Identifier: GPL-3.0-or-later
% Copyright (C) 2026 Alessandro Greco
%% Русское альтернативное описание EPUB.
\altfig{Конфигурация общедоступного onion-сервиса версии 3: HiddenServiceDir /var/lib/tor/my\_service/ и HiddenServicePort 80 127.0.0.1:8080. Tor создаёт в каталоге hostname, hs\_ed25519\_public\_key и hs\_ed25519\_secret\_key. Секретный ключ нужно защищать: злоумышленник, получивший его, может выдать себя за сервис. Адрес v3 содержит 56 символов base32, кодирующих открытый ключ Ed25519, контрольную сумму и байт версии; [v3].onion на схеме — лишь заполнитель. Соединения через Tor на виртуальный порт 80 перенаправляются веб-серверу, прослушивающему только 127.0.0.1:8080. Для общедоступного сервиса файл alice.auth не требуется; авторизация клиентов настраивается отдельно.}
% Публичный onion-сервис v3: конфигурация, ключи и локальное перенаправление
\begin{tikzpicture}[font=\sffamily\scriptsize,
  file/.style={draw,rounded corners=1pt,fill=torgrigio!8,minimum height=0.55cm,minimum width=3.4cm,font=\ttfamily\scriptsize},
  box/.style={draw,rounded corners=2pt,align=center,minimum height=0.8cm}]
  \node[box,fill=torviola!10,text width=4.6cm,align=left,font=\ttfamily\scriptsize] (conf) at (0,1.6)
    {HiddenServiceDir /var/lib/tor/my\_service/\\HiddenServicePort 80 127.0.0.1:8080};
  \node[font=\sffamily\tiny,text=torgrigio,above=1pt of conf] {torrc};
  \node[file] (h) at (0,0.1) {hostname};
  \node[file] (pk) at (0,-0.6) {hs\_ed25519\_public\_key};
  \node[file] (sk) at (0,-1.3) {hs\_ed25519\_secret\_key};
  \node[font=\sffamily\tiny,text=torgrigio,align=center] at (0,-1.95) {защитите каталог и закрытый ключ};
  \draw[->] (conf) -- (h);
  \node[box,fill=torverde!12,minimum width=2.4cm] (srv) at (6.4,-0.6) {веб-сервер\\127.0.0.1:8080};
  \node[box,fill=torviola!15,minimum width=2.4cm] (onion) at (6.4,1.2) {\texttt{[v3].onion}\\порт 80};
  \draw[->,torviola] (onion) -- node[right,font=\sffamily\tiny,align=left] {локальное\\перенаправление} (srv);
  \node[font=\sffamily\tiny,text=torgrigio] at (6.4,2.2) {v3: 56 base32};
  \draw[dashed,torgrigio,->] (pk.east) to[bend left=12] (onion.west);
  \node[anchor=west,font=\sffamily\tiny,text=avviso,align=left] at (4.4,-1.9) {похищенный секретный ключ\\позволяет подменить сервис};
  \draw[avviso,->] (sk.east) to[bend right=10] (4.3,-1.9);
\end{tikzpicture}

\caption{Директивы \opt{HiddenServiceDir} и \opt{HiddenServicePort} создают общедоступный onion-сервис v3. Tor генерирует в каталоге \file{hostname} и ключи Ed25519; адрес содержит 56 символов base32, полученных из открытого ключа, контрольной суммы и байта версии (в схеме указан заполнитель). Виртуальный порт 80 перенаправляется на \texttt{127.0.0.1:8080}. Кража закрытого ключа позволяет подменить сервис; для этого публичного примера файл \file{alice.auth} не нужен.}
\label{fig:servizio}
\end{figure}

\subsection{Как не раскрыть сервер}

Протокол скрывает адрес сервера, но небрежная конфигурация может его выдать. Типичные ошибки:
\begin{itemize}
  \item \textbf{Одновременная внешняя доступность сервиса.} Если веб-сервер слушает все интерфейсы, а не только \texttt{127.0.0.1}, сканирование Интернета может обнаружить его и сопоставить содержимое с onion-сервисом. Привязывайте сервер к localhost или, предпочтительно, используйте сокет Unix (\texttt{HiddenServicePort 80 unix:/path/socket}).
  \item \textbf{Утечки в ответах.} Страницы ошибок могут раскрывать имя сервера, HTTP-заголовки — внутренние адреса, TLS-сертификаты могут быть выпущены для публичного домена, а страницы состояния веб-сервера — оставлены доступными.
  \item \textbf{Внешние ресурсы.} Страницы, загружающие изображения или скрипты с обычных веб-сайтов, отправляют эти запросы за пределы onion-сервиса. В Tor Browser они всё равно проходят через Tor, поэтому IP-адрес посетителя не раскрывается, но сквозное шифрование onion-сервиса больше не покрывает эти ресурсы, и выходное реле может наблюдать метаданные соединения и незашифрованный HTTP-трафик (содержимое HTTPS по-прежнему защищено). Серверные приложения, самостоятельно обращающиеся к внешним сервисам за обновлениями или аналитикой, могут раскрыть адрес сервера.
  \item \textbf{Часы и временные характеристики.} Сервер, отображающий локальное время, связанное с конкретным часовым поясом, или становящийся недоступным одновременно с известным публичным сервером во время обслуживания, оставляет потенциально коррелируемые следы.
\end{itemize}
В \cref{cap:28-opsec} мы рассмотрим реальные случаи, где подобные ошибки помогли определить местоположение сервисов.

% ---------------------------------------------------------------------------
\section{Onion-Location}\label{sec:onionp-location}

Многие публичные сайты, включая крупные новостные организации, социальные сети и Tor Project, предлагают версию \texttt{.onion}. Обычный веб-сайт может сообщить посетителю Tor Browser о такой версии через HTTP-заголовок:
\begin{lstlisting}[numbers=none]
Onion-Location: http://<address>.onion/path
\end{lstlisting}
Эквивалентный элемент \texttt{meta} можно использовать в самой странице. Tor Browser после этого показывает кнопку в адресной строке, предлагающую перейти на onion-версию. Посетитель избегает выходное реле и получает соединение, аутентифицированное самим адресом; сайт принимает Tor-трафик без выхода через выходное реле, которые иногда блокируются или считаются подозрительными.

% ---------------------------------------------------------------------------
\section{Ограниченное обнаружение}\label{sec:onionp-ristretta}

Личный сервис, например SSH-доступ к домашнему серверу или административная панель, не обязан быть доступным всем, кто знает его адрес. Ограниченное обнаружение (restricted discovery, \cref{sec:os-ristretta}) шифрует второй уровень дескриптора так, чтобы расшифровать его могли только авторизованные клиенты. Конфигурация использует пары ключей X25519~\cite{torspec}:
\begin{description}
  \item[На стороне сервиса] Поместите по одному файлу \file{.auth} на клиента в каталог \file{authorized\_clients/} сервиса. Каждый файл содержит строку \texttt{descriptor:\allowbreak x25519:}, за которой следует открытый ключ клиента в base32.
  \item[На стороне клиента] Опция \opt{ClientOnionAuthDir} указывает каталог, содержащий по одному файлу \file{.auth\_private} на сервис. В каждом файле указывается адрес без \texttt{.onion}, затем \texttt{:\allowbreak descriptor:\allowbreak x25519:} и закрытый ключ клиента в base32.
\end{description}
Tor Browser позволяет ввести закрытый ключ при посещении сервиса, который его требует.

\textbf{Предупреждение о безопасности лаборатории:} скрипт \file{labs/ru/лаборатория/restricted\_discovery.py} во время начальной очистки удаляет существующие каталоги \file{restricted\_service/} и \file{onion\_auth/}. Запускайте его \emph{только} в выделенной одноразовой сети Chutney, убедившись, что в каталогах нет нужных данных. Не используйте его для настоящих onion-сервисов. При ручном выполнении сначала создайте файлы авторизации и только затем перезагрузите затронутые узлы.

\begin{laboratorio}[13]{Onion-сервис с ограниченным обнаружением}
\textbf{Цель.} Настроить второй onion-сервис на \texttt{test009h}, доступный только авторизованному клиенту, и проверить доступ с ключом и без него.\\
\textbf{Предварительные условия.} Лабораторная сеть должна быть запущена; локальный веб-сервер должен работать на порту 8080.\\
\textbf{Оценочное время.} 1 час.
\end{laboratorio}

\paragraph{Шаг 1: пара ключей клиента.} Сгенерируем её с помощью библиотеки \texttt{cryptography}:
Следующий самостоятельный пример Python печатает строку с открытым ключом авторизации. Закрытый ключ \texttt{priv} необходимо надёжно сохранить отдельно для создания файла \file{.auth\_private}; никому его не передавайте.
\begin{lstlisting}[language=Python]
import base64
from cryptography.hazmat.primitives.asymmetric.x25519 import X25519PrivateKey
from cryptography.hazmat.primitives.serialization import (
    Encoding, PrivateFormat, PublicFormat, NoEncryption
)

sk = X25519PrivateKey.generate()
priv = sk.private_bytes(Encoding.Raw, PrivateFormat.Raw, NoEncryption())
pub = sk.public_key().public_bytes(Encoding.Raw, PublicFormat.Raw)
b32 = lambda b: base64.b32encode(b).decode().rstrip("=")
print("descriptor:x25519:" + b32(pub))
\end{lstlisting}
Значения base32 должны записываться без завершающих символов заполнения \texttt{=}.
Фрагмент печатает только открытую строку авторизации и \emph{не сохраняет} закрытый ключ. При отдельном запуске сценария переменная \texttt{priv} исчезнет после завершения процесса. Для ручного выполнения заранее сохраните ключ в защищённый файл с ограниченными правами либо используйте автоматизированную лабораторную работу: она создаёт клиентский файл авторизации и удаляет временные ключи по окончании. Не следует смешивать ручной и автоматический сценарии.

\paragraph{Шаг 2: сервис.} Создайте \file{restricted\_service/authorized\_clients/alice.auth} с открытым ключом. Добавьте в \file{torrc} узла \texttt{test009h}:
\begin{lstlisting}[language=torrc,numbers=none]
HiddenServiceDir /lab/net/nodes/009h/restricted_service
HiddenServicePort 80 127.0.0.1:8080
\end{lstlisting}
Затем отправьте \texttt{HUP}. tor сгенерирует адрес и опубликует дескрипторы.

Следующая команда считывает настоящий onion-адрес этого экземпляра из файла \file{hostname}, который tor создаёт вместе с сервисом: у каждой сети свой адрес. Ошибка SOCKS и сообщения журнала ниже взяты из предыдущей эталонной сессии и могут отличаться при новом запуске.

\paragraph{Шаг 3: без ключа.} На клиенте сети:
\begin{lstlisting}[numbers=none]
$ ONION_ADDR=$(cat labs/ru/лаборатория/net/nodes/009h/restricted_service/hostname)
$ curl --noproxy "" --socks5-hostname 127.0.0.1:9008 \
    "http://${ONION_ADDR}/index.html"
(curl: SOCKS proxy error 97)
\end{lstlisting}
Журнал клиента сообщает:
\begin{lstlisting}[numbers=none]
[notice] Fail to decrypt descriptor for requested onion address.
         It is likely requiring client authorization.
\end{lstlisting}
Клиент нашёл и загрузил дескриптор, но не может расшифровать его второй уровень. Он даже не знает точек введения (introduction points) и потому не способен попытаться установить соединение.

\begin{attenzione}[Пустой каталог делает сервис недоступным]
Если \file{authorized\_clients/} существует, но не содержит ни одного файла, сервис в исходном эксперименте попытался опубликовать дескриптор, который серверы HSDir отклонили:
\begin{lstlisting}[numbers=none]
[warn] Uploading hidden service descriptor: http status 400 ("Invalid HS
       descriptor. Rejected.") response from dirserver 127.0.0.1:5106.
\end{lstlisting}
Соответствующая спецификация требует как минимум одну строку \texttt{auth-client} в дополнительно зашифрованном слое (superencrypted)~\spec{rend-spec}{Hidden service descriptors: encryption format}. Если ограниченное обнаружение активировано, но авторизованных клиентов нет, наблюдаемая конфигурация не создаёт ни одной такой строки, поэтому дескриптор оказывается недопустимым. Это эксплуатационная ловушка при отзыве последнего клиента или создании каталога до добавления ключей. Прерванный лабораторный сценарий также может оставить тестовую конфигурацию в незавершённом состоянии.
\end{attenzione}

\paragraph{Шаг 4: с ключом.} Запишите файл \file{.auth\_private} в каталог, указанный через \opt{ClientOnionAuthDir}, перезагрузите конфигурацию клиента и повторите запрос. Теперь страница будет доставлена. Скрипт \file{labs/ru/лаборатория/restricted\_discovery.py} автоматизирует эксперимент, очистку и восстановление \file{torrc}; его фактический запуск не проводился в рамках редакционной проверки.

% ---------------------------------------------------------------------------
\section{Другие требования}\label{sec:onionp-altre}

\subsection{Адреса с узнаваемым префиксом}

Адрес, начинающийся с узнаваемого слова, может помочь пользователю визуально его распознавать. Для этого специализированная программа вроде \emph{mkp224o} генерирует случайные ключи, пока адрес не начнётся с нужного префикса. Каждый дополнительный символ base32 увеличивает ожидаемое число попыток в 32 раза. Префикс из шести или семи символов можно получить за разумное время на обычном компьютере, но каждый следующий символ требует во много раз больше работы. Как обсуждалось в \cref{sec:os-indirizzo}, атакующий тоже может создать адрес с тем же префиксом. Префикс помогает распознаванию, но не аутентифицирует сервис. Обращайтесь с ключами, созданными сторонними программами, так же осторожно, как с любыми другими секретными ключами.

\subsection{Одношаговые сервисы}

Публичный веб-сайт, которому не требуется скрывать собственное местоположение, например onion-версия известной газеты, может пожертвовать анонимностью сервера ради скорости. При включённых \opt{HiddenServiceSingleHopMode} и \opt{HiddenServiceNonAnonymousMode} сервис строит одношаговые цепи к точкам введения (introduction points) и встречи (rendezvous points). Клиент остаётся анонимным как обычно. Публикация дескрипторов всё равно использует трёхузловые цепи: как мы выяснили в \cref{sec:lab-onion}, HSDir отклоняют прямые запросы, и это прямо отмечено в руководстве tor. Слово \enquote{NonAnonymous} в названии опции выбрано намеренно: экземпляр tor, настроенный таким образом, нельзя одновременно использовать как анонимный клиент. Кроме того, такое решение трудно обратить. После работы сервиса в неанонимном режиме его IP-адрес может быть связан с адресом \texttt{.onion}. Простого отключения опции недостаточно для восстановления анонимности. Необходимы новый адрес, новые ключи, новый \opt{HiddenServiceDir} и чистый экземпляр tor.

\subsection{Onionbalance}

Один процесс tor одновременно является узким местом производительности и единой точкой отказа. Onionbalance~\cite{onionbalance} позволяет нескольким серверам обслуживать один адрес. Каждый сервер запускает собственный onion-сервис со своими точками введения. Onionbalance хранит ключ основного адреса и публикует дескриптор, объединяющий точки введения всех серверов. Клиенты распределяются между ними прозрачно. Tor Project использует Onionbalance для собственных сайтов.

% ---------------------------------------------------------------------------
\section*{Итоги}
\begin{itemize}
  \item Onion-сервису достаточно двух строк конфигурации. Его секретный ключ \emph{контролирует} адрес и должен быть защищён и резервно сохранён.
  \item Сервер должен слушать только localhost или сокет Unix и не должен раскрывать сведения о себе через ответы или внешние соединения.
  \item Onion-Location сообщает посетителям Tor Browser о наличии onion-версии сайта.
  \item Ограниченное обнаружение использует ключи X25519: без авторизации клиент не может даже расшифровать дескриптор.
  \item Адреса с узнаваемым префиксом, одношаговые сервисы и Onionbalance решают задачи узнаваемости, скорости и масштабирования.
\end{itemize}

% ---------------------------------------------------------------------------
\section*{Упражнения}

\begin{esercizio}[1]\label{es:21-vanity}
Сколько попыток в среднем требуется для генерации шестисимвольного префикса base32? А восьмисимвольного? Если компьютер проверяет десять миллионов ключей в секунду, сколько времени займёт каждый вариант?
\end{esercizio}

\begin{esercizio}[2]\label{es:21-fuga}
Onion-сервис показывает на страницах ошибок строку \texttt{Apache/2.4.62 (Debian) Server at 192.0.2.10 Port 80}. Какую информацию она раскрывает? Как избежать этой утечки?
\end{esercizio}

\begin{esercizio}[2]\label{es:21-revoca}
Как в Лабораторной работе~13 отозвать доступ Алисы? Что произойдёт с клиентами, которые уже загрузили дескриптор?
\end{esercizio}

\begin{esercizio}[3]\label{es:21-socket}
Измените лабораторный сервис так, чтобы соединения перенаправлялись в Unix socket, а не на локальный TCP-порт. Какие дополнительные преимущества это даёт помимо отсутствия открытого порта?
\end{esercizio}

% ---------------------------------------------------------------------------
\section*{Дополнительная литература}
\begin{itemize}
  \item Разделы руководства tor для \opt{HiddenServiceDir}, \opt{HiddenServicePort} и client authorization.
  \item Документация Tor Project по onion-сервисам и Onionbalance~\cite{onionbalance}.
\end{itemize}
